% frontmatter/actividades.tex -- las actividades, una a una, para el
% adulto: cada estación trae las suyas (ver notes/01-plan.md), y esta
% página tiene que seguir cabiendo en una (tools/check_pages.py comprueba
% que el día 1 está en la página 5).
\section*{Las actividades}

Cada actividad llega el día que se usa por primera vez, y después vuelve
de vez en cuando. Las instrucciones de cada página, en letra pequeña y
gris, las lee un adulto, si hace falta.

{\small
\begin{multicols}{2}
\raggedright
\subsection*{Desde el otoño}
\begin{itemize}[leftmargin=5mm, itemsep=1pt, topsep=0pt]
\item \textbf{\lblClasifica}: leer cada tarjeta y unirla con su caja
  -- por ejemplo, \emph{lo necesito} o \emph{lo quiero}.
\item \textbf{\lblUne}: unir cada cosa de la izquierda con la suya, a
  la derecha.
\item \textbf{\lblVF}: leer cada frase y rodear la V si es verdad, o la
  F si es mentira.
\item \textbf{\lblTrueque}: si tantas cosas valen tantas otras, ¿cuántas
  te dan por estas? Se puede dibujar.
\item \textbf{\lblDibuja}: dibujar lo que dice.
\item \textbf{\lblRepasa} (los viernes): la palabra de la semana, lo que
  ya sabe hacer, para marcarlo, y un dibujo. Cada diez páginas hay un
  pequeño cartel de celebración.
\end{itemize}
\end{multicols}}
\newpage
